% Options for packages loaded elsewhere
% Options for packages loaded elsewhere
\PassOptionsToPackage{unicode}{hyperref}
\PassOptionsToPackage{hyphens}{url}
\PassOptionsToPackage{dvipsnames,svgnames,x11names}{xcolor}
%
\documentclass[
  letterpaper,
  DIV=11,
  numbers=noendperiod]{scrartcl}
\usepackage{xcolor}
\usepackage{amsmath,amssymb}
\setcounter{secnumdepth}{5}
\usepackage{iftex}
\ifPDFTeX
  \usepackage[T1]{fontenc}
  \usepackage[utf8]{inputenc}
  \usepackage{textcomp} % provide euro and other symbols
\else % if luatex or xetex
  \usepackage{unicode-math} % this also loads fontspec
  \defaultfontfeatures{Scale=MatchLowercase}
  \defaultfontfeatures[\rmfamily]{Ligatures=TeX,Scale=1}
\fi
\usepackage{lmodern}
\ifPDFTeX\else
  % xetex/luatex font selection
\fi
% Use upquote if available, for straight quotes in verbatim environments
\IfFileExists{upquote.sty}{\usepackage{upquote}}{}
\IfFileExists{microtype.sty}{% use microtype if available
  \usepackage[]{microtype}
  \UseMicrotypeSet[protrusion]{basicmath} % disable protrusion for tt fonts
}{}
\makeatletter
\@ifundefined{KOMAClassName}{% if non-KOMA class
  \IfFileExists{parskip.sty}{%
    \usepackage{parskip}
  }{% else
    \setlength{\parindent}{0pt}
    \setlength{\parskip}{6pt plus 2pt minus 1pt}}
}{% if KOMA class
  \KOMAoptions{parskip=half}}
\makeatother
% Make \paragraph and \subparagraph free-standing
\makeatletter
\ifx\paragraph\undefined\else
  \let\oldparagraph\paragraph
  \renewcommand{\paragraph}{
    \@ifstar
      \xxxParagraphStar
      \xxxParagraphNoStar
  }
  \newcommand{\xxxParagraphStar}[1]{\oldparagraph*{#1}\mbox{}}
  \newcommand{\xxxParagraphNoStar}[1]{\oldparagraph{#1}\mbox{}}
\fi
\ifx\subparagraph\undefined\else
  \let\oldsubparagraph\subparagraph
  \renewcommand{\subparagraph}{
    \@ifstar
      \xxxSubParagraphStar
      \xxxSubParagraphNoStar
  }
  \newcommand{\xxxSubParagraphStar}[1]{\oldsubparagraph*{#1}\mbox{}}
  \newcommand{\xxxSubParagraphNoStar}[1]{\oldsubparagraph{#1}\mbox{}}
\fi
\makeatother

\usepackage{color}
\usepackage{fancyvrb}
\newcommand{\VerbBar}{|}
\newcommand{\VERB}{\Verb[commandchars=\\\{\}]}
\DefineVerbatimEnvironment{Highlighting}{Verbatim}{commandchars=\\\{\}}
% Add ',fontsize=\small' for more characters per line
\usepackage{framed}
\definecolor{shadecolor}{RGB}{241,243,245}
\newenvironment{Shaded}{\begin{snugshade}}{\end{snugshade}}
\newcommand{\AlertTok}[1]{\textcolor[rgb]{0.68,0.00,0.00}{#1}}
\newcommand{\AnnotationTok}[1]{\textcolor[rgb]{0.37,0.37,0.37}{#1}}
\newcommand{\AttributeTok}[1]{\textcolor[rgb]{0.40,0.45,0.13}{#1}}
\newcommand{\BaseNTok}[1]{\textcolor[rgb]{0.68,0.00,0.00}{#1}}
\newcommand{\BuiltInTok}[1]{\textcolor[rgb]{0.00,0.23,0.31}{#1}}
\newcommand{\CharTok}[1]{\textcolor[rgb]{0.13,0.47,0.30}{#1}}
\newcommand{\CommentTok}[1]{\textcolor[rgb]{0.37,0.37,0.37}{#1}}
\newcommand{\CommentVarTok}[1]{\textcolor[rgb]{0.37,0.37,0.37}{\textit{#1}}}
\newcommand{\ConstantTok}[1]{\textcolor[rgb]{0.56,0.35,0.01}{#1}}
\newcommand{\ControlFlowTok}[1]{\textcolor[rgb]{0.00,0.23,0.31}{\textbf{#1}}}
\newcommand{\DataTypeTok}[1]{\textcolor[rgb]{0.68,0.00,0.00}{#1}}
\newcommand{\DecValTok}[1]{\textcolor[rgb]{0.68,0.00,0.00}{#1}}
\newcommand{\DocumentationTok}[1]{\textcolor[rgb]{0.37,0.37,0.37}{\textit{#1}}}
\newcommand{\ErrorTok}[1]{\textcolor[rgb]{0.68,0.00,0.00}{#1}}
\newcommand{\ExtensionTok}[1]{\textcolor[rgb]{0.00,0.23,0.31}{#1}}
\newcommand{\FloatTok}[1]{\textcolor[rgb]{0.68,0.00,0.00}{#1}}
\newcommand{\FunctionTok}[1]{\textcolor[rgb]{0.28,0.35,0.67}{#1}}
\newcommand{\ImportTok}[1]{\textcolor[rgb]{0.00,0.46,0.62}{#1}}
\newcommand{\InformationTok}[1]{\textcolor[rgb]{0.37,0.37,0.37}{#1}}
\newcommand{\KeywordTok}[1]{\textcolor[rgb]{0.00,0.23,0.31}{\textbf{#1}}}
\newcommand{\NormalTok}[1]{\textcolor[rgb]{0.00,0.23,0.31}{#1}}
\newcommand{\OperatorTok}[1]{\textcolor[rgb]{0.37,0.37,0.37}{#1}}
\newcommand{\OtherTok}[1]{\textcolor[rgb]{0.00,0.23,0.31}{#1}}
\newcommand{\PreprocessorTok}[1]{\textcolor[rgb]{0.68,0.00,0.00}{#1}}
\newcommand{\RegionMarkerTok}[1]{\textcolor[rgb]{0.00,0.23,0.31}{#1}}
\newcommand{\SpecialCharTok}[1]{\textcolor[rgb]{0.37,0.37,0.37}{#1}}
\newcommand{\SpecialStringTok}[1]{\textcolor[rgb]{0.13,0.47,0.30}{#1}}
\newcommand{\StringTok}[1]{\textcolor[rgb]{0.13,0.47,0.30}{#1}}
\newcommand{\VariableTok}[1]{\textcolor[rgb]{0.07,0.07,0.07}{#1}}
\newcommand{\VerbatimStringTok}[1]{\textcolor[rgb]{0.13,0.47,0.30}{#1}}
\newcommand{\WarningTok}[1]{\textcolor[rgb]{0.37,0.37,0.37}{\textit{#1}}}

\usepackage{longtable,booktabs,array}
\newcounter{none} % for unnumbered tables
\usepackage{calc} % for calculating minipage widths
% Correct order of tables after \paragraph or \subparagraph
\usepackage{etoolbox}
\makeatletter
\patchcmd\longtable{\par}{\if@noskipsec\mbox{}\fi\par}{}{}
\makeatother
% Allow footnotes in longtable head/foot
\IfFileExists{footnotehyper.sty}{\usepackage{footnotehyper}}{\usepackage{footnote}}
\makesavenoteenv{longtable}
\usepackage{graphicx}
\makeatletter
\newsavebox\pandoc@box
\newcommand*\pandocbounded[1]{% scales image to fit in text height/width
  \sbox\pandoc@box{#1}%
  \Gscale@div\@tempa{\textheight}{\dimexpr\ht\pandoc@box+\dp\pandoc@box\relax}%
  \Gscale@div\@tempb{\linewidth}{\wd\pandoc@box}%
  \ifdim\@tempb\p@<\@tempa\p@\let\@tempa\@tempb\fi% select the smaller of both
  \ifdim\@tempa\p@<\p@\scalebox{\@tempa}{\usebox\pandoc@box}%
  \else\usebox{\pandoc@box}%
  \fi%
}
% Set default figure placement to htbp
\def\fps@figure{htbp}
\makeatother





\setlength{\emergencystretch}{3em} % prevent overfull lines

\providecommand{\tightlist}{%
  \setlength{\itemsep}{0pt}\setlength{\parskip}{0pt}}



 


\KOMAoption{captions}{tableheading}
\makeatletter
\@ifpackageloaded{caption}{}{\usepackage{caption}}
\AtBeginDocument{%
\ifdefined\contentsname
  \renewcommand*\contentsname{Table of contents}
\else
  \newcommand\contentsname{Table of contents}
\fi
\ifdefined\listfigurename
  \renewcommand*\listfigurename{List of Figures}
\else
  \newcommand\listfigurename{List of Figures}
\fi
\ifdefined\listtablename
  \renewcommand*\listtablename{List of Tables}
\else
  \newcommand\listtablename{List of Tables}
\fi
\ifdefined\figurename
  \renewcommand*\figurename{Figure}
\else
  \newcommand\figurename{Figure}
\fi
\ifdefined\tablename
  \renewcommand*\tablename{Table}
\else
  \newcommand\tablename{Table}
\fi
}
\@ifpackageloaded{float}{}{\usepackage{float}}
\floatstyle{ruled}
\@ifundefined{c@chapter}{\newfloat{codelisting}{h}{lop}}{\newfloat{codelisting}{h}{lop}[chapter]}
\floatname{codelisting}{Listing}
\newcommand*\listoflistings{\listof{codelisting}{List of Listings}}
\makeatother
\makeatletter
\makeatother
\makeatletter
\@ifpackageloaded{caption}{}{\usepackage{caption}}
\@ifpackageloaded{subcaption}{}{\usepackage{subcaption}}
\makeatother
\usepackage{bookmark}
\IfFileExists{xurl.sty}{\usepackage{xurl}}{} % add URL line breaks if available
\urlstyle{same}
\hypersetup{
  pdftitle={GII Index Reduction Analysis Report},
  pdfauthor={GII Analysis Team},
  colorlinks=true,
  linkcolor={blue},
  filecolor={Maroon},
  citecolor={Blue},
  urlcolor={Blue},
  pdfcreator={LaTeX via pandoc}}


\title{GII Index Reduction Analysis Report}
\author{GII Analysis Team}
\date{2026-06-04}
\begin{document}
\maketitle

\renewcommand*\contentsname{Table of contents}
{
\setcounter{tocdepth}{3}
\tableofcontents
}

\section{Executive Summary}\label{executive-summary}

This report presents the findings of the analysis to reduce the number
of indicators in the Global Innovation Index (GII) while maintaining the
discriminatory power of the index. Based on our analysis, we recommend a
subset of 15 indicators that capture approximately 86\% of the
variability of the full 50-indicator index.

\section{Introduction}\label{introduction}

When creating indicator dashboards and composite indices, indicator
selection often involves narrowing down from an initially larger set.
Besides data quality and relevance considerations, statistical methods
are essential to assess whether indicators within a well-defined
component are sufficiently aligned. A key statistical question is: how
much of the information content (``variation'') of the initial subset is
preserved in a potential subset?

This analysis applies two methodological approaches to the GII data:

\begin{enumerate}
\def\labelenumi{\arabic{enumi}.}
\tightlist
\item
  \textbf{Country Perspective}: Assessing how similar countries are
  using all indicators versus a subset, using distance matrices and
  clustering techniques
\item
  \textbf{Variable Perspective}: Using Principal Component Analysis
  (PCA) to measure how much information a subset captures compared to
  the full set
\end{enumerate}

\section{Methodology}\label{methodology}

\subsection{Statistical Methods
Explanation}\label{statistical-methods-explanation}

To determine how much of the information content of the full indicator
set is preserved in a potential subset, we use a subset selection
approach based on Principal Component Analysis (PCA). Rather than simply
ranking individual variables by importance, this method evaluates entire
subsets of variables according to how well they reproduce the PCA
structure of the full dataset.

The approach uses a criterion called \(r_m\) that measures how much of
the overall multivariate information in the full PCA can be recovered
from a given subset of variables. This is a more coherent objective for
dimension reduction than individual variable selection, as it identifies
combinations of variables that best preserve the complete data
structure.

A key advantage of this method is its built-in protection against
redundancy. For example, if multiple indicators are highly correlated
(nearly identical), including one in the subset provides most of the
information, and adding the others contributes relatively little
additional ability to reconstruct the principal components.

The mathematical formula for the selection criterion is:

\[ r_m = \sqrt{ \frac{\sum_{i=1}^{p}\lambda_i r_{mi}^2} {\sum_{j=1}^{p}\lambda_j} }\]

where: - \(\lambda_i\) is the eigenvalue of principal component \(i\)
from the full dataset (representing the variance explained by that
component) - \(m\) indexes a specific candidate subset - \(r_{mi}^2\) is
the coefficient of determination (\(R^2\)) between principal component
\(i\) and the candidate subset of variables. Since \(PC_{i}\) is a
linear combination of the original variables, the \(R^2\) of the subset
measures how well the principle component can be reconstructed from the
subset.

Conceptually, this formula computes a weighted average of how well the
subset predicts each principal component, with weights given by the
variance explained by each component, then normalizes by the total
variance. The resulting \(r_m^2\) value can be interpreted as the
percentage of total variation explained by the subset, providing an
intuitive measure of information retention.

This method assumes that ``information'' is defined by the variance
structure captured by PCA, which is appropriate for statistical data
reduction. However, it may downweight variables with low variance but
high policy relevance, which is an important consideration for
substantive index development.

To prevent the method from over-emphasizing statistical data-reduction,
we also implement a methodology where sub-sets of variables are selected
by stratifying the sampling across all pillars of the GII index.

\subsection{Sampling Implementation}\label{sampling-implementation}

To ensure representative and statistically robust variable selection
across the conceptual framework of the Global Innovation Index, we
implement a stratified sampling approach that integrates both
proportional allocation and clustering principles. This methodology,
inspired by survey sampling techniques across geographic regions,
ensures balanced representation while maintaining information content.

Our sampling strategy consists of the following components:

\begin{enumerate}
\def\labelenumi{\arabic{enumi}.}
\item
  \textbf{Stratification by Conceptual Pillars}: The seven GII pillars
  (IN.1--IN.5, OUT.6, OUT.7) serve as strata, ensuring conceptual
  coverage in every subset.
\item
  \textbf{Proportional Allocation}: For each subset of target size, we
  allocate variables to pillars proportional to their size in the
  indicator framework. This initial allocation ensures larger pillars
  contribute more variables while maintaining minimum representation.
\item
  \textbf{Allocation Adjustment}: To achieve exact subset sizes, we
  implement a dynamic adjustment algorithm that incrementally assigns
  additional slots to the largest pillars until the target size is
  reached, handling rounding discrepancies from proportional allocation.
\item
  \textbf{Hierarchical Clustering within Pillars}: Within each pillar,
  we apply hierarchical clustering based on correlation structure to
  identify distinct indicator clusters. This step ensures diversity by
  preventing selection of redundant, highly correlated indicators.
\item
  \textbf{Cluster-Based Sampling}: We sample one indicator from each
  identified cluster within a pillar when possible, ensuring internal
  diversity. This clustering is performed using Ward's method on
  dissimilarity matrices derived from correlation coefficients.
\item
  \textbf{Fallback Mechanism}: When cluster-based sampling cannot fill
  the required allocation (due to insufficient clusters or data
  limitations), we supplement with random sampling from remaining
  indicators in the pillar.
\end{enumerate}

This multi-stage approach combines the statistical rigor of stratified
sampling with the information-preserving properties of clustering,
resulting in subsets that are both conceptually balanced and
statistically representative of the full indicator space. By
incorporating cluster analysis within pillars, our method goes beyond
simple random stratified sampling to actively avoid redundancy while
preserving the multidimensional structure of the GII framework.

\subsection{Data Description}\label{data-description}

\begin{itemize}
\tightlist
\item
  here description of GII 2025 data. 139 regions across 78 indicators,
  we impute the missing values using k-NN.
\end{itemize}

\section{Key Findings}\label{key-findings}

\subsection{Distance Matrix Analysis}\label{distance-matrix-analysis}

The analysis compared distance matrices for different subset sizes
against the full dataset:

\begin{Shaded}
\begin{Highlighting}[]
\CommentTok{\# Read and display subset size comparison results}
\NormalTok{results }\OtherTok{\textless{}{-}} \FunctionTok{read.csv}\NormalTok{(}\StringTok{"outputs/subset\_size\_comparison.csv"}\NormalTok{)}
\NormalTok{knitr}\SpecialCharTok{::}\FunctionTok{kable}\NormalTok{(results, }\AttributeTok{caption =} \StringTok{"Distance Matrix Comparison by Subset Size"}\NormalTok{, }\AttributeTok{digits =} \DecValTok{2}\NormalTok{)}
\end{Highlighting}
\end{Shaded}

\begin{longtable}[]{@{}rrr@{}}
\caption{Distance Matrix Comparison by Subset Size}\tabularnewline
\toprule\noalign{}
size & mean\_correlation & max\_correlation \\
\midrule\noalign{}
\endfirsthead
\toprule\noalign{}
size & mean\_correlation & max\_correlation \\
\midrule\noalign{}
\endhead
\bottomrule\noalign{}
\endlastfoot
10 & 0.71 & 0.80 \\
15 & 0.78 & 0.86 \\
20 & 0.83 & 0.88 \\
25 & 0.86 & 0.90 \\
30 & 0.89 & 0.93 \\
\end{longtable}

\textbf{Finding:} 15 indicators achieve 0.781 mean correlation with full
dataset, significantly better than random subsets (p \textless{} 0.001).

\subsection{Clustering Results}\label{clustering-results}

\begin{itemize}
\tightlist
\item
  \textbf{Full dataset (FTD):} 14 significant clusters
\item
  \textbf{Reduced dataset (RTD):} 11 significant clusters\\
\item
  \textbf{Rand Index:} 95.7\%
\item
  \textbf{Most significant loss:} Luxembourg, Finland, Sweden, Denmark,
  Netherlands (3 clusters in FTD → 1 in RTD)
\end{itemize}

\textbf{Finding:} The reduced 44-indicator set produces a coarser but
generally consistent clustering structure.

\subsection{PCA Analysis}\label{pca-analysis}

The table below shows the number of principal components needed to
capture specified levels of variance:

\begin{Shaded}
\begin{Highlighting}[]
\CommentTok{\# Read and display PCA components}
\NormalTok{pca\_components }\OtherTok{\textless{}{-}} \FunctionTok{read.csv}\NormalTok{(}\StringTok{"outputs/pca\_components.csv"}\NormalTok{)}
\NormalTok{knitr}\SpecialCharTok{::}\FunctionTok{kable}\NormalTok{(pca\_components, }\AttributeTok{caption =} \StringTok{"Components Needed for Target Variance Levels"}\NormalTok{)}
\end{Highlighting}
\end{Shaded}

\begin{longtable}[]{@{}lr@{}}
\caption{Components Needed for Target Variance Levels}\tabularnewline
\toprule\noalign{}
Variance\_Level & Components\_Needed \\
\midrule\noalign{}
\endfirsthead
\toprule\noalign{}
Variance\_Level & Components\_Needed \\
\midrule\noalign{}
\endhead
\bottomrule\noalign{}
\endlastfoot
80\% & 17 \\
90\% & 26 \\
95\% & 36 \\
\end{longtable}

\textbf{Finding:} First 17 components capture 80\% of total variance,
suggesting substantial dimensionality reduction potential.

\subsection{Variable Perspective
Results}\label{variable-perspective-results}

\begin{Shaded}
\begin{Highlighting}[]
\CommentTok{\# Read and display r\_m results}
\NormalTok{r\_m\_results }\OtherTok{\textless{}{-}} \FunctionTok{read.csv}\NormalTok{(}\StringTok{"outputs/r\_m\_subset\_size\_comparison.csv"}\NormalTok{)}
\NormalTok{knitr}\SpecialCharTok{::}\FunctionTok{kable}\NormalTok{(r\_m\_results, }\AttributeTok{format =} \StringTok{"simple"}\NormalTok{, }\AttributeTok{digits =} \DecValTok{2}\NormalTok{)}
\end{Highlighting}
\end{Shaded}

{\def\LTcaptype{none} % do not increment counter
\begin{longtable}[]{@{}rrr@{}}
\toprule\noalign{}
size & mean\_r\_m & max\_r\_m \\
\midrule\noalign{}
\endhead
\bottomrule\noalign{}
\endlastfoot
10 & 0.74 & 0.76 \\
15 & 0.79 & 0.81 \\
20 & 0.85 & 0.86 \\
25 & 0.87 & 0.89 \\
30 & 0.88 & 0.90 \\
35 & 0.89 & 0.91 \\
40 & 0.90 & 0.92 \\
\end{longtable}
}

\textbf{Finding:} Subsets consistently outperform random selections,
with 15-variable subset achieving 0.792 of total variability.

\subsection{Recommended Variable Subset (Top
15)}\label{recommended-variable-subset-top-15}

{\def\LTcaptype{none} % do not increment counter
\begin{longtable}[]{@{}lll@{}}
\toprule\noalign{}
Rank & Variable Name & Domain \\
\midrule\noalign{}
\endhead
\bottomrule\noalign{}
\endlastfoot
1 & Rule of Law & Institutions \\
2 & Government Effectiveness & Wellbeing Today \\
3 & Regulatory Quality & Institutions \\
4 & Female Employment Gap & Inclusiveness \\
5 & Economic Complexity & Societal Resilience \\
6 & Public Services (Digital) & Societal Resilience \\
7 & ICT Access (5G) & Resources for Future \\
8 & Life Expectancy & Resources for Future \\
9 & University Res. Coll. & Societal Resilience \\
10 & Science Quality & Resources for Future \\
11 & TLDs Mix & Connectivity \\
12 & GH Rec. Sent. (No Bot) & Innovation \\
13 & Demographic Dividend & Societal Resilience \\
14 & Intellectual Property Protection & Resources for Future \\
15 & Documentation Quality & Business \\
\end{longtable}
}

\textbf{Finding:} This subset achieves r\_m = 0.827, capturing 86\% of
total variability.

\subsection{Variable Perspective Results (by
Pillar)}\label{variable-perspective-results-by-pillar}

To ensure representation across the conceptual pillars of the Global
Innovation Index, we implemented a pillar-constrained variable selection
approach. This method guarantees that each random subset includes at
least one indicator from each of the seven pillars (IN.1--IN.5, OUT.6,
OUT.7) before filling remaining slots with random indicators.

{\def\LTcaptype{none} % do not increment counter
\begin{longtable}[]{@{}llll@{}}
\toprule\noalign{}
Subset Size & Mean r\_m & Max r\_m & Performance vs.~Random \\
\midrule\noalign{}
\endhead
\bottomrule\noalign{}
\endlastfoot
10 & 0.731 & 0.757 & \textgreater99.9th percentile \\
15 & 0.794 & 0.815 & \textgreater99.9th percentile \\
20 & 0.843 & 0.862 & \textgreater99.9th percentile \\
25 & 0.879 & 0.893 & \textgreater99.9th percentile \\
30 & 0.907 & 0.921 & \textgreater99.9th percentile \\
\end{longtable}
}

\textbf{Finding:} The pillar-constrained approach ensures balanced
representation across GII domains while achieving 79.4\% mean
variability capture with the 15-variable subset.

\subsection{Optimal 15-Variable Subset (by
Pillar)}\label{optimal-15-variable-subset-by-pillar}

The optimal 15-variable subset with pillar representation:

{\def\LTcaptype{none} % do not increment counter
\begin{longtable}[]{@{}
  >{\raggedright\arraybackslash}p{(\linewidth - 4\tabcolsep) * \real{0.2069}}
  >{\raggedright\arraybackslash}p{(\linewidth - 4\tabcolsep) * \real{0.5172}}
  >{\raggedright\arraybackslash}p{(\linewidth - 4\tabcolsep) * \real{0.2759}}@{}}
\toprule\noalign{}
\begin{minipage}[b]{\linewidth}\raggedright
Rank
\end{minipage} & \begin{minipage}[b]{\linewidth}\raggedright
Variable Name
\end{minipage} & \begin{minipage}[b]{\linewidth}\raggedright
Domain
\end{minipage} \\
\midrule\noalign{}
\endhead
\bottomrule\noalign{}
\endlastfoot
1 & Government Effectiveness & Wellbeing Today \\
2 & Researchers, Full-Time Equivalent per 1,000 Labour Force & Resources
for Future \\
3 & Logistics Performance Index & Connectivity \\
4 & Financial Support to SMEs & Resources for Future \\
5 & Researchers, Tertiary Sector & Societal Resilience \\
6 & Economic Complexity & Societal Resilience \\
7 & Entertainment Media & Societal Resilience \\
8 & Regulatory Quality & Institutions \\
9 & Rule of Law & Institutions \\
10 & PISA Score & Resources for Future \\
11 & Research and Development Expenditure as Percentage of GDP &
Resources for Future \\
12 & Knowledge-intensive Employment & Societal Resilience \\
13 & GERD Financing from Business & Resources for Future \\
14 & Public-Private Publication Co-authorships & Societal Resilience \\
15 & THE University Ranking Score & Resources for Future \\
\end{longtable}
}

\textbf{Finding:} This pillar-representative subset achieves r\_m =
0.794, capturing 79.4\% of total variability while ensuring coverage
across all conceptual domains.

\section{Comparison: 15 vs.~44
Indicators}\label{comparison-15-vs.-44-indicators}

{\def\LTcaptype{none} % do not increment counter
\begin{longtable}[]{@{}llll@{}}
\toprule\noalign{}
Metric & 15 Indicators & 44 Indicators & Difference \\
\midrule\noalign{}
\endhead
\bottomrule\noalign{}
\endlastfoot
Mean Correlation & 0.834 & 0.851 & -1.7\% \\
Max Correlation & 0.898 & 0.938 & -4.0\% \\
Rand Index & 0.827 & 0.908 & -8.1\% \\
Variables & 15 & 44 & -66\% \\
\end{longtable}
}

\textbf{Finding:} The 15-indicator subset loses only 1.7\% in mean
correlation compared to the 44-indicator set, while reducing indicator
count by 66\%.

\section{Conclusions}\label{conclusions}

\begin{enumerate}
\def\labelenumi{\arabic{enumi}.}
\item
  \textbf{Significant reduction possible:} A 15-indicator subset
  captures 86\% of the variability of the full 50-indicator index
\item
  \textbf{Statistical significance:} The 15-indicator subset
  significantly outperforms random subsets (p \textless{} 0.001)
\item
  \textbf{Consistent performance:} Results are stable across different
  analysis methods (distance matrices, PCA, r\_m metric)
\item
  \textbf{Trade-off:} The 15-indicator subset loses some clustering
  specificity (Rand Index: 82.7\% vs.~95.7\%) but maintains strong
  overall correlation (0.834)
\item
  \textbf{Policy implications:} A 66\% reduction in indicators could
  simplify data collection and interpretation while maintaining most of
  the index's discriminatory power
\end{enumerate}

\section{Limitations}\label{limitations}

Users should understand what information is lost with reduced indicator
sets 1. Analysis focused on one year (2025); temporal stability needs
assessment 2. Geographic coverage limited to available data (75
countries) 3. Methods assume linear relationships; non-linear effects
may exist 4. Domain-specific validity not tested

\section{Data Availability}\label{data-availability}

All data and code are available at: - Processed data:
\texttt{data/gii\_data\_standardized.csv} - Analysis scripts:
\texttt{scripts/} - Results: \texttt{outputs/}

\section{References}\label{references}

\begin{enumerate}
\def\labelenumi{\arabic{enumi}.}
\tightlist
\item
  European Commission/JRC. (2023). Sustainable and Inclusive Well-being
  Index: Methodology Report.
\item
  UN HLEG. (2024). Beyond GDP Dashboard: Conceptual Framework.
\item
  Global Innovation Index 2025: Data and analysis.
\end{enumerate}

\section{Appendix}\label{appendix}

\section{Data Imputation
Documentation}\label{data-imputation-documentation}

Missing data is an inherent challenge in composite index construction.
Our analysis employed k-Nearest Neighbors (kNN) imputation to address
missing values before conducting the main analyses. This section
documents the imputation process and its results.

\subsection{Imputation Process}\label{imputation-process}

\begin{itemize}
\tightlist
\item
  \textbf{Method}: k-Nearest Neighbors (kNN) imputation
\item
  \textbf{Parameters}: k=5 neighbors, rowmax=0.9, colmax=0.8
\item
  \textbf{Implementation}: R \texttt{VIM} package
\item
  \textbf{Indicators Used}: Level 1 indicators from IMETA metadata
\end{itemize}

\subsection{Missing Data Statistics}\label{missing-data-statistics}

The table below summarizes the imputation process:

\begin{Shaded}
\begin{Highlighting}[]
\CommentTok{\# Read and display imputation report}
\NormalTok{imputation\_report }\OtherTok{\textless{}{-}} \FunctionTok{read.csv}\NormalTok{(}\StringTok{"data/imputation\_report.csv"}\NormalTok{)}
\NormalTok{knitr}\SpecialCharTok{::}\FunctionTok{kable}\NormalTok{(imputation\_report, }\AttributeTok{caption =} \StringTok{"Imputation Process Summary"}\NormalTok{)}
\end{Highlighting}
\end{Shaded}

\begin{longtable}[]{@{}lr@{}}
\caption{Imputation Process Summary}\tabularnewline
\toprule\noalign{}
Metric & Count \\
\midrule\noalign{}
\endfirsthead
\toprule\noalign{}
Metric & Count \\
\midrule\noalign{}
\endhead
\bottomrule\noalign{}
\endlastfoot
Total Missing Values & 1211 \\
Variables with Missing Data & 58 \\
Countries with Missing Data & 135 \\
\end{longtable}

\section{References on Imputation}\label{references-on-imputation}

The kNN imputation method preserves the data structure while filling
missing values based on similar observations. This approach is
particularly suitable for composite indices where maintaining the
relationships between indicators is crucial for valid comparisons.

\textbf{Method}: k-Nearest Neighbors (kNN) imputation with K = 5

\section{Descriptive Statistics
Tables}\label{descriptive-statistics-tables}

\subsection{Descriptive Statistics - Original
Data}\label{descriptive-statistics---original-data}

The table below presents the descriptive statistics for all Level 1
indicators in the original dataset before imputation:

\begin{Shaded}
\begin{Highlighting}[]
\CommentTok{\# Read and display original descriptive statistics}
\NormalTok{original\_stats }\OtherTok{\textless{}{-}} \FunctionTok{read.csv}\NormalTok{(}\StringTok{"data/descriptive\_stats\_original.csv"}\NormalTok{)}
\NormalTok{knitr}\SpecialCharTok{::}\FunctionTok{kable}\NormalTok{(original\_stats, }\AttributeTok{caption =} \StringTok{"Descriptive Statistics {-} Original Data (n=78 indicators)"}\NormalTok{, }\AttributeTok{digits =} \DecValTok{2}\NormalTok{)}
\end{Highlighting}
\end{Shaded}

\begin{longtable}[]{@{}
  >{\raggedright\arraybackslash}p{(\linewidth - 12\tabcolsep) * \real{0.2237}}
  >{\raggedleft\arraybackslash}p{(\linewidth - 12\tabcolsep) * \real{0.1447}}
  >{\raggedleft\arraybackslash}p{(\linewidth - 12\tabcolsep) * \real{0.1316}}
  >{\raggedleft\arraybackslash}p{(\linewidth - 12\tabcolsep) * \real{0.1053}}
  >{\raggedleft\arraybackslash}p{(\linewidth - 12\tabcolsep) * \real{0.1579}}
  >{\raggedleft\arraybackslash}p{(\linewidth - 12\tabcolsep) * \real{0.1184}}
  >{\raggedleft\arraybackslash}p{(\linewidth - 12\tabcolsep) * \real{0.1184}}@{}}
\caption{Descriptive Statistics - Original Data (n=78
indicators)}\tabularnewline
\toprule\noalign{}
\begin{minipage}[b]{\linewidth}\raggedright
Indicator
\end{minipage} & \begin{minipage}[b]{\linewidth}\raggedleft
Mean
\end{minipage} & \begin{minipage}[b]{\linewidth}\raggedleft
Median
\end{minipage} & \begin{minipage}[b]{\linewidth}\raggedleft
Min
\end{minipage} & \begin{minipage}[b]{\linewidth}\raggedleft
Max
\end{minipage} & \begin{minipage}[b]{\linewidth}\raggedleft
Kurtosis
\end{minipage} & \begin{minipage}[b]{\linewidth}\raggedleft
Skewness
\end{minipage} \\
\midrule\noalign{}
\endfirsthead
\toprule\noalign{}
\begin{minipage}[b]{\linewidth}\raggedright
Indicator
\end{minipage} & \begin{minipage}[b]{\linewidth}\raggedleft
Mean
\end{minipage} & \begin{minipage}[b]{\linewidth}\raggedleft
Median
\end{minipage} & \begin{minipage}[b]{\linewidth}\raggedleft
Min
\end{minipage} & \begin{minipage}[b]{\linewidth}\raggedleft
Max
\end{minipage} & \begin{minipage}[b]{\linewidth}\raggedleft
Kurtosis
\end{minipage} & \begin{minipage}[b]{\linewidth}\raggedleft
Skewness
\end{minipage} \\
\midrule\noalign{}
\endhead
\bottomrule\noalign{}
\endlastfoot
PolStab & 2.23 & 2.15 & 0.70 & 4.45 & -0.06 & 0.56 \\
GovEff & 0.15 & 0.03 & -1.75 & 2.32 & -0.76 & 0.29 \\
RegQua & 0.16 & 0.11 & -2.03 & 2.31 & -0.77 & 0.20 \\
RuleOL & 0.07 & -0.12 & -2.15 & 1.97 & -0.81 & 0.28 \\
PolicyDB & 4.03 & 3.98 & 1.83 & 6.45 & -0.75 & 0.02 \\
EntreprenPCplus & 4.72 & 4.72 & 2.92 & 7.25 & -0.32 & 0.44 \\
ExpEdu & 4.25 & 4.16 & 0.35 & 9.39 & 0.68 & 0.28 \\
GovFundPup & 19.60 & 19.41 & 6.00 & 51.06 & 1.33 & 0.72 \\
SchLifeExp & 14.17 & 14.36 & 6.34 & 20.85 & -0.31 & -0.28 \\
PISA & 439.44 & 439.29 & 337.45 & 579.03 & -0.91 & 0.06 \\
PupTeachR & 15.89 & 13.53 & 6.14 & 68.13 & 9.23 & 2.41 \\
TertEnrol & 51.51 & 54.40 & 2.71 & 166.67 & -0.11 & 0.36 \\
SciEngGrad & 23.29 & 23.14 & 9.56 & 41.10 & -0.41 & 0.17 \\
TertMob & 7.34 & 4.24 & 0.00 & 69.73 & 14.40 & 3.26 \\
ResPPop & 2269.70 & 922.47 & 4.50 & 10413.00 & 0.39 & 1.18 \\
RDExp & 1.01 & 0.61 & 0.01 & 6.35 & 4.13 & 1.87 \\
RDCompExp & 926.35 & 0.00 & 0.00 & 36142.79 & 57.81 & 6.95 \\
QSRank & 21.29 & 10.43 & 0.00 & 97.63 & 0.19 & 1.12 \\
ICTAccWIPO5G20 & 83.02 & 88.25 & 24.33 & 100.00 & 0.42 & -1.09 \\
ICTUseWIPO2 & 79.65 & 84.47 & 29.84 & 99.36 & 1.60 & -1.45 \\
GovOLServ & 0.68 & 0.72 & 0.17 & 1.00 & -0.98 & -0.39 \\
ElecOut & 4936.33 & 3183.45 & 29.76 & 53263.16 & 22.90 & 3.97 \\
LogPerf & 3.12 & 3.10 & 2.10 & 4.30 & -1.13 & 0.19 \\
GroCap & 24.49 & 23.70 & 11.60 & 44.56 & 0.51 & 0.72 \\
GDPEner & 12.23 & 11.06 & 2.16 & 41.21 & 3.70 & 1.56 \\
LowCarbonEner & 22.00 & 18.38 & 0.01 & 82.64 & 0.69 & 1.05 \\
ISOEnv & 2.05 & 0.99 & 0.04 & 11.87 & 2.72 & 1.78 \\
FinanceSSGEMplus & 4.49 & 4.49 & 1.99 & 6.90 & -0.57 & 0.09 \\
DomCred2Priv & 60.33 & 48.53 & 8.14 & 248.76 & 2.30 & 1.48 \\
MicroLoanIMFFAS & 2.09 & 1.05 & 0.00 & 17.34 & 11.33 & 3.18 \\
MarkCap & 80.85 & 37.45 & 0.20 & 1507.77 & 51.10 & 6.72 \\
VenCapInvPB & 0.65 & 0.14 & 0.01 & 29.23 & 95.18 & 9.53 \\
VenCapRecPB & 0.24 & 0.08 & 0.01 & 5.50 & 63.34 & 7.25 \\
LateVCInt & 0.32 & 0.03 & 0.00 & 11.43 & 52.16 & 6.93 \\
JointVenPB & 0.21 & 0.06 & 0.00 & 4.55 & 52.79 & 6.49 \\
AppTariffWTO & 4.16 & 2.34 & 0.00 & 47.54 & 38.87 & 4.98 \\
HHI & 0.18 & 0.14 & 0.07 & 0.55 & 2.78 & 1.57 \\
DomMarkSc & 1370.77 & 253.12 & 4.14 & 37072.09 & 45.03 & 6.41 \\
Demogdividend & 39.00 & 35.41 & 17.63 & 66.87 & -1.03 & 0.52 \\
KnowIntsEmp & 27.07 & 23.94 & 2.02 & 68.01 & -0.99 & 0.30 \\
FemEmpDeg & 13.49 & 12.78 & 0.14 & 31.04 & -1.17 & 0.25 \\
GERDPerf & 0.77 & 0.38 & 0.00 & 5.90 & 6.95 & 2.25 \\
GERDFinBus & 34.31 & 38.44 & 0.00 & 80.84 & -1.12 & -0.07 \\
PublicPrivCopubs & 2.01 & 1.34 & 0.22 & 9.07 & 2.73 & 1.74 \\
UniResColl & 3.97 & 3.86 & 2.30 & 6.72 & -0.20 & 0.38 \\
THEunirank & 55.74 & 49.40 & 26.70 & 97.10 & -1.04 & 0.56 \\
StatClustDev & 4.00 & 3.90 & 2.59 & 5.40 & -0.83 & 0.24 \\
PatFam & 1.30 & 0.03 & 0.00 & 73.67 & 108.33 & 10.07 \\
IntelPropPay & 1.05 & 0.54 & 0.00 & 23.25 & 61.02 & 7.02 \\
HiTechImp & 9.33 & 8.01 & 1.30 & 55.19 & 20.27 & 3.57 \\
ICTSevImp & 1.75 & 1.44 & 0.01 & 16.01 & 29.45 & 4.34 \\
FDINetInfl & 3.93 & 2.50 & -161.34 & 224.25 & 63.06 & 3.20 \\
ResTalBus & 33.25 & 32.27 & 0.30 & 81.91 & -1.23 & 0.20 \\
PatOrig & 2.54 & 0.62 & 0.00 & 61.56 & 37.33 & 5.79 \\
PCTInventors & 0.62 & 0.07 & 0.00 & 7.57 & 11.16 & 3.22 \\
UtilModOrig & 1.55 & 0.16 & 0.00 & 88.51 & 77.80 & 8.85 \\
SciTechArt & 13.42 & 9.94 & 0.55 & 47.72 & 0.40 & 1.11 \\
CitDoc & 466.22 & 281.00 & 53.00 & 3213.00 & 7.50 & 2.33 \\
GrowPPP & 0.83 & 0.86 & -4.57 & 6.26 & 0.71 & -0.01 \\
UnicornVal & 1.09 & 0.00 & 0.00 & 21.85 & 26.08 & 4.65 \\
CompSoftSpend & 0.24 & 0.21 & 0.00 & 1.26 & 2.84 & 1.39 \\
HiTechManuf & 25.19 & 21.99 & 1.04 & 79.65 & 0.03 & 0.73 \\
IntelPropRec & 0.48 & 0.07 & 0.00 & 5.36 & 8.62 & 2.95 \\
EcComplexity & 0.09 & -0.05 & -2.18 & 2.29 & -0.58 & 0.22 \\
HTExpHKG & 4.57 & 1.51 & 0.02 & 48.39 & 11.70 & 3.01 \\
ICTServExp & 2.99 & 1.68 & 0.00 & 36.98 & 25.52 & 4.45 \\
ISOQual & 5.46 & 3.23 & 0.04 & 29.87 & 3.16 & 1.76 \\
CorpIAs & 41.66 & 56.72 & -371.30 & 89.79 & 26.36 & -4.63 \\
TradmkOrig & 36.72 & 28.95 & 0.00 & 202.33 & 6.90 & 2.19 \\
BrandVal & 3.59 & 0.78 & 0.00 & 21.12 & 1.58 & 1.60 \\
IndustDes & 2.21 & 0.92 & 0.00 & 23.28 & 11.16 & 3.17 \\
CultServExp & 0.96 & 0.46 & 0.00 & 18.31 & 50.72 & 6.32 \\
FilmsOmdia & 5.01 & 2.67 & 0.00 & 46.93 & 17.44 & 3.59 \\
EntMedia & 0.93 & 0.47 & 0.02 & 3.59 & -0.24 & 0.87 \\
CreaGoodExp & 1.27 & 0.38 & 0.00 & 11.50 & 6.56 & 2.55 \\
TLDsMix & 117.16 & 22.31 & 0.12 & 2538.57 & 43.80 & 5.80 \\
GHRecSentNoBot & 1326116.27 & 410927.87 & 2571.55 & 14402260.97 & 13.66
& 2.93 \\
AppCrea & 1379627.73 & 181351.44 & 0.04 & 98744342.48 & 117.05 &
10.73 \\
\end{longtable}

\subsection{Descriptive Statistics - Complete
Data}\label{descriptive-statistics---complete-data}

The table below presents the descriptive statistics for all Level 1
indicators in the complete dataset after kNN imputation and case-wise
cleaning:

\begin{Shaded}
\begin{Highlighting}[]
\CommentTok{\# Read and display complete descriptive statistics}
\NormalTok{complete\_stats }\OtherTok{\textless{}{-}} \FunctionTok{read.csv}\NormalTok{(}\StringTok{"data/descriptive\_stats\_complete.csv"}\NormalTok{)}
\NormalTok{knitr}\SpecialCharTok{::}\FunctionTok{kable}\NormalTok{(complete\_stats, }\AttributeTok{caption =} \StringTok{"Descriptive Statistics {-} Complete Data (n=78 indicators)"}\NormalTok{, }\AttributeTok{digits =} \DecValTok{2}\NormalTok{)}
\end{Highlighting}
\end{Shaded}

\begin{longtable}[]{@{}
  >{\raggedright\arraybackslash}p{(\linewidth - 12\tabcolsep) * \real{0.2237}}
  >{\raggedleft\arraybackslash}p{(\linewidth - 12\tabcolsep) * \real{0.1447}}
  >{\raggedleft\arraybackslash}p{(\linewidth - 12\tabcolsep) * \real{0.1316}}
  >{\raggedleft\arraybackslash}p{(\linewidth - 12\tabcolsep) * \real{0.1053}}
  >{\raggedleft\arraybackslash}p{(\linewidth - 12\tabcolsep) * \real{0.1579}}
  >{\raggedleft\arraybackslash}p{(\linewidth - 12\tabcolsep) * \real{0.1184}}
  >{\raggedleft\arraybackslash}p{(\linewidth - 12\tabcolsep) * \real{0.1184}}@{}}
\caption{Descriptive Statistics - Complete Data (n=78
indicators)}\tabularnewline
\toprule\noalign{}
\begin{minipage}[b]{\linewidth}\raggedright
Indicator
\end{minipage} & \begin{minipage}[b]{\linewidth}\raggedleft
Mean
\end{minipage} & \begin{minipage}[b]{\linewidth}\raggedleft
Median
\end{minipage} & \begin{minipage}[b]{\linewidth}\raggedleft
Min
\end{minipage} & \begin{minipage}[b]{\linewidth}\raggedleft
Max
\end{minipage} & \begin{minipage}[b]{\linewidth}\raggedleft
Kurtosis
\end{minipage} & \begin{minipage}[b]{\linewidth}\raggedleft
Skewness
\end{minipage} \\
\midrule\noalign{}
\endfirsthead
\toprule\noalign{}
\begin{minipage}[b]{\linewidth}\raggedright
Indicator
\end{minipage} & \begin{minipage}[b]{\linewidth}\raggedleft
Mean
\end{minipage} & \begin{minipage}[b]{\linewidth}\raggedleft
Median
\end{minipage} & \begin{minipage}[b]{\linewidth}\raggedleft
Min
\end{minipage} & \begin{minipage}[b]{\linewidth}\raggedleft
Max
\end{minipage} & \begin{minipage}[b]{\linewidth}\raggedleft
Kurtosis
\end{minipage} & \begin{minipage}[b]{\linewidth}\raggedleft
Skewness
\end{minipage} \\
\midrule\noalign{}
\endhead
\bottomrule\noalign{}
\endlastfoot
PolStab & 2.23 & 2.15 & 0.70 & 4.45 & -0.06 & 0.56 \\
GovEff & 0.15 & 0.03 & -1.75 & 2.32 & -0.76 & 0.29 \\
RegQua & 0.16 & 0.11 & -2.03 & 2.31 & -0.77 & 0.20 \\
RuleOL & 0.07 & -0.12 & -2.15 & 1.97 & -0.81 & 0.28 \\
PolicyDB & 4.03 & 3.98 & 1.83 & 6.45 & -0.70 & 0.02 \\
EntreprenPCplus & 4.63 & 4.54 & 2.92 & 7.25 & 0.18 & 0.63 \\
ExpEdu & 4.24 & 4.14 & 0.35 & 9.39 & 0.70 & 0.30 \\
GovFundPup & 18.48 & 17.96 & 6.00 & 51.06 & 1.78 & 0.83 \\
SchLifeExp & 13.92 & 13.94 & 6.34 & 20.85 & -0.49 & -0.12 \\
PISA & 419.11 & 403.60 & 337.45 & 579.03 & -0.50 & 0.68 \\
PupTeachR & 15.80 & 13.28 & 6.14 & 68.13 & 9.43 & 2.43 \\
TertEnrol & 50.58 & 53.42 & 2.71 & 166.67 & -0.09 & 0.40 \\
SciEngGrad & 22.88 & 22.88 & 9.56 & 41.10 & -0.38 & 0.25 \\
TertMob & 6.61 & 3.14 & 0.00 & 69.73 & 16.80 & 3.50 \\
ResPPop & 1922.18 & 669.77 & 4.50 & 10413.00 & 0.94 & 1.40 \\
RDExp & 0.88 & 0.42 & 0.01 & 6.35 & 5.34 & 2.11 \\
RDCompExp & 933.55 & 0.00 & 0.00 & 36142.79 & 58.14 & 6.97 \\
QSRank & 21.29 & 10.43 & 0.00 & 97.63 & 0.19 & 1.12 \\
ICTAccWIPO5G20 & 82.88 & 88.19 & 24.33 & 100.00 & 0.36 & -1.07 \\
ICTUseWIPO2 & 78.81 & 83.81 & 29.84 & 99.36 & 1.13 & -1.33 \\
GovOLServ & 0.68 & 0.73 & 0.17 & 1.00 & -0.98 & -0.40 \\
ElecOut & 4663.61 & 2879.19 & 29.76 & 53263.16 & 24.24 & 4.08 \\
LogPerf & 3.04 & 2.90 & 2.10 & 4.30 & -1.00 & 0.44 \\
GroCap & 24.48 & 23.70 & 11.60 & 44.56 & 0.55 & 0.74 \\
GDPEner & 12.00 & 10.93 & 2.16 & 41.21 & 4.16 & 1.65 \\
LowCarbonEner & 22.00 & 18.38 & 0.01 & 82.64 & 0.69 & 1.05 \\
ISOEnv & 2.05 & 0.99 & 0.04 & 11.87 & 2.72 & 1.78 \\
FinanceSSGEMplus & 4.25 & 4.11 & 1.99 & 6.90 & -0.51 & 0.47 \\
DomCred2Priv & 60.08 & 48.30 & 8.14 & 248.76 & 2.27 & 1.47 \\
MicroLoanIMFFAS & 1.43 & 0.63 & 0.00 & 17.34 & 24.87 & 4.49 \\
MarkCap & 58.57 & 25.34 & 0.20 & 1507.77 & 84.32 & 8.54 \\
VenCapInvPB & 0.56 & 0.11 & 0.01 & 29.23 & 114.84 & 10.45 \\
VenCapRecPB & 0.22 & 0.07 & 0.01 & 5.50 & 70.18 & 7.61 \\
LateVCInt & 0.25 & 0.01 & 0.00 & 11.43 & 68.34 & 7.91 \\
JointVenPB & 0.18 & 0.04 & 0.00 & 4.55 & 62.78 & 7.06 \\
AppTariffWTO & 4.16 & 2.34 & 0.00 & 47.54 & 38.87 & 4.98 \\
HHI & 0.19 & 0.17 & 0.07 & 0.55 & 1.87 & 1.17 \\
DomMarkSc & 1370.77 & 253.12 & 4.14 & 37072.09 & 45.03 & 6.41 \\
Demogdividend & 39.00 & 35.41 & 17.63 & 66.87 & -1.03 & 0.52 \\
KnowIntsEmp & 26.67 & 23.31 & 2.02 & 68.01 & -0.93 & 0.37 \\
FemEmpDeg & 13.06 & 11.20 & 0.14 & 31.04 & -1.08 & 0.34 \\
GERDPerf & 0.50 & 0.06 & 0.00 & 5.90 & 11.27 & 2.90 \\
GERDFinBus & 28.45 & 28.73 & 0.00 & 80.84 & -1.05 & 0.34 \\
PublicPrivCopubs & 2.01 & 1.34 & 0.22 & 9.07 & 2.73 & 1.74 \\
UniResColl & 3.95 & 3.83 & 2.30 & 6.72 & -0.13 & 0.42 \\
THEunirank & 51.65 & 42.20 & 26.70 & 97.10 & -0.32 & 1.00 \\
StatClustDev & 3.99 & 3.84 & 2.59 & 5.40 & -0.80 & 0.28 \\
PatFam & 1.29 & 0.03 & 0.00 & 73.67 & 109.14 & 10.11 \\
IntelPropPay & 1.05 & 0.54 & 0.00 & 23.25 & 61.02 & 7.02 \\
HiTechImp & 9.32 & 8.03 & 1.30 & 55.19 & 20.45 & 3.59 \\
ICTSevImp & 1.75 & 1.44 & 0.01 & 16.01 & 29.45 & 4.34 \\
FDINetInfl & 3.91 & 2.49 & -161.34 & 224.25 & 63.54 & 3.22 \\
ResTalBus & 23.49 & 11.34 & 0.30 & 81.91 & -0.67 & 0.85 \\
PatOrig & 2.51 & 0.61 & 0.00 & 61.56 & 37.91 & 5.84 \\
PCTInventors & 0.58 & 0.06 & 0.00 & 7.57 & 12.48 & 3.38 \\
UtilModOrig & 1.05 & 0.16 & 0.00 & 88.51 & 129.76 & 11.39 \\
SciTechArt & 13.42 & 9.94 & 0.55 & 47.72 & 0.40 & 1.11 \\
CitDoc & 466.22 & 281.00 & 53.00 & 3213.00 & 7.50 & 2.33 \\
GrowPPP & 0.83 & 0.86 & -4.57 & 6.26 & 0.74 & 0.00 \\
UnicornVal & 1.09 & 0.00 & 0.00 & 21.85 & 26.08 & 4.65 \\
CompSoftSpend & 0.24 & 0.20 & 0.00 & 1.26 & 2.85 & 1.39 \\
HiTechManuf & 22.20 & 18.35 & 1.04 & 79.65 & 0.26 & 0.92 \\
IntelPropRec & 0.48 & 0.07 & 0.00 & 5.36 & 8.62 & 2.95 \\
EcComplexity & 0.08 & -0.06 & -2.18 & 2.29 & -0.50 & 0.25 \\
HTExpHKG & 4.53 & 1.51 & 0.02 & 48.39 & 11.79 & 3.02 \\
ICTServExp & 2.99 & 1.68 & 0.00 & 36.98 & 25.52 & 4.45 \\
ISOQual & 5.46 & 3.23 & 0.04 & 29.87 & 3.16 & 1.76 \\
CorpIAs & 32.37 & 39.35 & -371.30 & 89.79 & 32.65 & -4.58 \\
TradmkOrig & 36.74 & 28.99 & 0.00 & 202.33 & 7.11 & 2.21 \\
BrandVal & 3.20 & 0.47 & 0.00 & 21.12 & 2.26 & 1.77 \\
IndustDes & 2.12 & 0.84 & 0.00 & 23.28 & 11.91 & 3.27 \\
CultServExp & 0.86 & 0.42 & 0.00 & 18.31 & 57.63 & 6.72 \\
FilmsOmdia & 4.47 & 2.69 & 0.00 & 46.93 & 24.78 & 4.16 \\
EntMedia & 0.54 & 0.17 & 0.02 & 3.59 & 2.52 & 1.80 \\
CreaGoodExp & 1.27 & 0.38 & 0.00 & 11.50 & 6.63 & 2.56 \\
TLDsMix & 117.16 & 22.31 & 0.12 & 2538.57 & 43.80 & 5.80 \\
GHRecSentNoBot & 1357449.63 & 411718.75 & 2571.55 & 14402260.97 & 13.14
& 2.84 \\
AppCrea & 1305696.24 & 155386.88 & 0.04 & 98744342.48 & 124.55 &
11.07 \\
\end{longtable}




\end{document}
